% Zdroj: PDF priklady-jaro-2024.pdf, fyzická str. 24. Značka: specializace UI-SU, UI-ZPJ. Zařazeno: S1 (informované prohledávání).
\section{Informované prohledávání}
\szzotazka{jaro 2024}{32}{Informované prohledávání (A*, heuristika, 8-puzzle)}

\noindent\textbf{Zadání.}\par
Uvažme hlavolam na obrázku níže. Jedná se o mřížku $3 \times 3$ s 8 označenými kostičkami a jedním volným místem. Naším cílem je ze zadaného počátečního stavu uspořádat kostičky do předem daného cílového stavu, ve kterém jsou kostičky uspořádány od 1 do 8 zleva doprava, shora dolů. V každém kroku můžeme posunout právě jednu kostičku na sousední volné místo. Navíc chceme provést minimální možný počet kroků.

\begin{center}\includegraphics[width=0.32\linewidth]{8-puzzle.png}\end{center}

\begin{enumerate}[nosep]
  \item Formalizujte zadaný problém jako informovaný prohledávací problém. Speciálně popište, jak vypadá stav, jaký je přechodový model a jaký je test cílového stavu.
  \item Pro vyřešení definovaného problému navrhněte optimální informovaný prohledávací algoritmus. Popište jeho průběh (například pomocí pseudokódu).
  \item Navrhněte heuristickou funkci pro výše zmíněný hlavolam. Jaké vlastnosti má vaše heuristická funkce? Zaručuje použití této heuristiky s algoritmem z předchozího bodu nalezení optimálního řešení? \textit{Nápověda: Napsali jste algoritmus jako tree search nebo graph search?}
\end{enumerate}

\medskip
\noindent\textbf{Náčrt řešení.}\par
\begin{enumerate}[nosep]
  \item Stav je rozložení kostiček, přechodový model je pohnutí jednou kostičkou, čímž dostanu jiné rozložení (nezapomenout na cenu akce $= 1$), kontrola cílového stavu zkontroluje, zda se rozložení shoduje se zadaným cílem.
  \item A*, je jedno jestli graph search nebo tree search.
  \item Například taxicab distance. U graph search musí být heuristika monotonní/konzistentní. U tree search stačí přípustnost/spodní odhad. Všechny rozumné heuristiky budou splňovat oboje.
\end{enumerate}

\medskip
\noindent\textit{Doplňující vysvětlení (nad rámec oficiálního náčrtu):} Dvě klasické heuristiky pro zadaný počáteční stav: $h_1$ = počet kostek mimo cílovou pozici ($h_1=6$, správně leží jen $3$ a~$7$) a~$h_2$ = součet manhattanských vzdáleností kostek od cíle ($h_2=1+1+0+2+2+1+0+2=9$ pro kostky $1$ až $8$). Obě jsou přípustné (jeden tah posune jednu kostku o~jedno pole, takže skutečný počet tahů nepodstřelí) i~konzistentní (tahem klesne $h$ nejvýš o~$1$, což je cena kroku); $h_2\ge h_1$, tedy $h_2$ dominuje. Přípustnost zaručuje optimalitu A* ve stromovém prohledávání, konzistence v~grafovém (s~closed listem). Poznámka k~instanci: počáteční stav má $11$ inverzí (liché číslo), cílový $0$, a~posun kostky paritu inverzí nemění. Zadaná instance proto do cíle \emph{není řešitelná}; A* (v~grafové verzi) by po prohledání všech $9!/2=181\,440$ dosažitelných stavů ohlásil neúspěch. Otázka ale konkrétní řešení nežádá.
